% 使用 XeLaTeX 编译；文件编码为 UTF-8。
\documentclass[UTF8,a4paper,12pt,fontset=fandol]{ctexart}

\usepackage[margin=2.5cm]{geometry}
\usepackage{amsmath,amssymb,amsthm}
\usepackage{hyperref}
\hypersetup{unicode=true,hidelinks,pdftitle={高等数学学习笔记}}

\newtheorem{theorem}{定理}[section]
\newtheorem{lemma}[theorem]{引理}
\newtheorem{corollary}[theorem]{推论}
\theoremstyle{definition}
\newtheorem{definition}{定义}[section]
\newtheorem{example}{例题}[section]
\theoremstyle{remark}
\newtheorem*{remark}{注}

\title{高等数学学习笔记}
\author{}
\date{}

\begin{document}
\maketitle

% 待提供本项目其他对话内容后，补充对应的定义、定理、推导和例题。
% 以下为章节结构示例，取消注释并替换内容即可使用。
% \section{章节标题}
% \subsection{基本概念}
% \begin{definition}
%   定义内容。
% \end{definition}
% \subsection{定理与推导}
% \begin{theorem}
%   定理内容。
% \end{theorem}
% \begin{proof}
%   证明过程。
% \end{proof}
% \subsection{典型例题}
% \begin{example}
%   题目内容。
% \end{example}
% \begin{proof}[解]
%   解题过程。
% \end{proof}

\end{document}
